\section{Rare observations and stochastic stability}\label{sec:sgd}

Unit marginal variance does not control the size of individual stochastic
updates. The following scalar calculation isolates this distinction using
known population centering and scaling. It is an elementary moment calculation
that motivates the multivariate problem.

\begin{proposition}[Scalar batch stability]\label{prop:sgd-stability}
Let $X\sim\operatorname{Bernoulli}(q)$ with $0<q<1$, and set
\[
 Z=\frac{X-q}{\sqrt{q(1-q)}},\qquad
 \mu_4=\E Z^4=\frac{(1-q)^3+q^3}{q(1-q)}.
\]
Fit the noiseless model $Y=Z\theta^*$ by SGD on squared loss
$(Z\theta-Y)^2/2$, with constant step size $\eta$ and batches of size
$B\ge1$. Each batch consists of fresh independent observations and is
independent of the previous iterates. For a deterministic initial error
$e_0=\theta_0-\theta^*$,
\[
 \E e_k^2=
 \left[1-2\eta+\eta^2\left\{1+\frac{\mu_4-1}{B}\right\}\right]^k e_0^2.
\]
For $e_0\ne0$, the expected squared error tends to zero precisely when
\begin{equation}\label{eq:sgd-stability}
 0<\eta<\frac{2}{1+(\mu_4-1)/B}.
\end{equation}
\end{proposition}
\begin{proof}
Write $H_k=B^{-1}\sum_{i=1}^B Z_{ki}^2$. The update gives
$e_{k+1}=(1-\eta H_k)e_k$.
Since $\E Z^2=1$, independence within each batch gives
$\E H_k=1$ and $\E H_k^2=1+(\mu_4-1)/B$.
Independence from the past yields the displayed recurrence for the second
moment. Its nonnegative multiplicative factor is less than one exactly
under \eqref{eq:sgd-stability}.
\end{proof}

The population Hessian equals one for every $q$, while $\mu_4\sim q^{-1}$
as $q\downarrow0$. If $B=B_q$ satisfies $B_q q\to0$, the stability bound is
asymptotic to $2B_q q$. If $B_q q\to\infty$, it tends to two, the stability
bound for the population gradient update. At $B_q q\to\gamma\in(0,\infty)$,
it tends to $2\gamma/(1+\gamma)$. These are mean-square stability statements
for the specified online noiseless model, not guarantees for Adam or for
minibatches drawn without replacement from a fixed sample.

This calculation also distinguishes mean convergence from mean-square
convergence: $\E e_k=(1-\eta)^k e_0$, which can converge while the second
moment grows. A covariance-only description of the population loss therefore
does not determine stochastic stability.

In one dimension, rescaling a feature can be absorbed into the learning rate.
The substantive normalization question requires several feature frequencies
and relative coordinate updates. With common and rare predictors together,
a shared learning rate must accommodate their different update moments.
Deriving that interaction with the transient dependence effect in
Section~\ref{sec:flow}, under label noise and a fixed observation budget, is
the next research objective.
